\documentclass[c,8pt]{beamer}

\setbeamercovered{dynamic}

% --- Colors ---
\definecolor{MyYellow}{RGB}{220,180,0}
\definecolor{MyAmber}{RGB}{250,229,33}
\definecolor{MyLightYellow}{RGB}{255,255,130}
\definecolor{MyBlue}{RGB}{0,102,204}
\definecolor{MyRed}{RGB}{204,0,0}
\definecolor{MyLightBlue}{RGB}{50,100,255}
\definecolor{MyRefBlue1}{RGB}{73,157,220}
\definecolor{MyRefBlue2}{RGB}{135,179,229}
\definecolor{MyGreen}{RGB}{0,153,76}
\definecolor{MyPurple}{RGB}{112,10,214}
\definecolor{MyOrange}{RGB}{255,111,31}
\definecolor{themeBlue}{rgb}{0.29,0.58,0.87}
\definecolor{amber}{rgb}{1.0,0.75,0.0}
\definecolor{greenrow}{RGB}{220,245,220}
\definecolor{yellowrow}{RGB}{255,245,200}
\definecolor{redrow}{RGB}{255,230,230}
\definecolor{bluerow}{RGB}{220,230,255}
\definecolor{bleucit}{rgb}{0.2,0.4,0.6}

\setbeamercolor{structure}{fg=themeBlue}

% --- Packages ---
\usepackage[english]{babel}
\usepackage{amsmath,mathrsfs,amsfonts,amssymb,bm}
\usepackage{graphicx,caption,subcaption}
\usepackage{soul,cancel}
\usepackage{multicol}
\usepackage{hyperref}
\usepackage{comment}
\usepackage{ifthen}
\usepackage{ragged2e}
\usepackage{booktabs,colortbl,array}
\usepackage{algorithm,algpseudocode}
\usepackage{epstopdf,pifont}
\usepackage{tikz}
\usetikzlibrary{shapes,arrows,calc,positioning,arrows.meta,
                decorations.pathmorphing,decorations.markings,
                matrix,mindmap,trees,automata}
\pgfdeclarelayer{background}
\pgfdeclarelayer{foreground}
\pgfsetlayers{background,main,foreground}
\usepackage[skins,theorems]{tcolorbox}
\usepackage{varwidth}

% --- Theme ---
\usetheme{Madrid}
\useinnertheme{default}
\usecolortheme{default}
\usefonttheme{default}
\useoutertheme[subsection=false]{smoothbars}

% --- Navigation ---
\beamertemplatenavigationsymbolsempty
\expandafter\def\expandafter\insertshorttitle\expandafter{RL Robotics 2026}
\setbeamercovered{transparent}
\setbeamercolor{footline}{fg=black}
\setbeamerfont{footline}{series=\bfseries}
\AtBeginSection[]{
  \begin{frame}
    \frametitle{Table of Contents}
    \tableofcontents[currentsection]
  \end{frame}
}

% --- Block color overrides ---
\newcommand\blueblock[2]{{\setbeamercolor{block title}{fg=white,bg=themeBlue}\begin{block}{#1}#2\end{block}}}
\newcommand\yellowblock[2]{{\setbeamercolor{block title}{fg=black,bg=MyAmber}\begin{block}{#1}#2\end{block}}}
\newcommand\redblock[2]{{\setbeamercolor{block title}{fg=white,bg=MyRed}\begin{block}{#1}#2\end{block}}}
\newcommand\greenblock[2]{{\setbeamercolor{block title}{fg=white,bg=MyGreen}\begin{block}{#1}#2\end{block}}}
\newcommand\purpleblock[2]{{\setbeamercolor{block title}{fg=white,bg=MyPurple}\begin{block}{#1}#2\end{block}}}

\AtBeginEnvironment{theorem}{\setbeamercolor{block title}{fg=white,bg=themeBlue}}
\AtBeginEnvironment{assumption}{\setbeamercolor{block title}{fg=white,bg=themeBlue}}
\AtBeginEnvironment{proposition}{\setbeamercolor{block title}{fg=white,bg=themeBlue}}
\AtBeginEnvironment{definition}{\setbeamercolor{block title}{fg=white,bg=MyGreen}}

% --- Theorem environments ---
\newtheorem{assumption}{Assumption}
\newtheorem{proposition}{Proposition}

% --- Math operators ---
\newcommand{\pderiv}[2]{\dfrac{\partial #1}{\partial #2}}
\def\RR{\mathbb{R}}
\def\EE{\mathbb{E}}
\def\PP{\mathbb{P}}
\DeclareMathOperator{\argmin}{argmin}
\DeclareMathOperator{\argmax}{argmax}
\DeclareMathOperator{\trace}{tr}
\newcommand{\dddt}{{\dfrac{{\rm d}}{{\rm d}t}}}
\newcommand{\ddt}{{\tfrac{{\rm d}}{{\rm d}t}}}
\renewcommand{\Re}{\mathfrak{Re}}
\renewcommand{\Im}{\mathfrak{Im}}
\newcommand{\btVFill}{\vskip0pt plus 1filll}
\def\beeq#1{\begin{equation}{#1}\end{equation}}
\newcommand\qmx[1]{\left(\begin{matrix}#1\end{matrix}\right)}
\newcommand\pmx[1]{\left(\begin{matrix}#1\end{matrix}\right)}
\newcommand\smallmat[1]{\left[\begin{smallmatrix}#1\end{smallmatrix}\right]}
\newcommand\pdef[1]{\mathbb{S}_{\succ0}^{#1}}
\newcommand\psemidef[1]{\mathbb{S}_{\succeq0}^{#1}}
\newcommand\wnorm[2]{\lvert #1\rvert^2_{#2}}
\newcommand\spec[1]{\text{spec}(#1)}

% --- Calligraphic shortcuts ---
\newcommand{\cA}{\mathcal{A}}  \newcommand{\cB}{\mathcal{B}}
\newcommand{\cC}{\mathcal{C}}  \newcommand{\cD}{\mathcal{D}}
\newcommand{\cE}{\mathcal{E}}  \newcommand{\cF}{\mathcal{F}}
\newcommand{\cG}{\mathcal{G}}  \newcommand{\cH}{\mathcal{H}}
\newcommand{\cI}{\mathcal{I}}  \newcommand{\cJ}{\mathcal{J}}
\newcommand{\cK}{\mathcal{K}}  \newcommand{\cL}{\mathcal{L}}
\newcommand{\cM}{\mathcal{M}}  \newcommand{\cN}{\mathcal{N}}
\newcommand{\cO}{\mathcal{O}}  \newcommand{\cP}{\mathcal{P}}
\newcommand{\cQ}{\mathcal{Q}}  \newcommand{\cR}{\mathcal{R}}
\newcommand{\cS}{\mathcal{S}}  \newcommand{\cT}{\mathcal{T}}
\newcommand{\cU}{\mathcal{U}}  \newcommand{\cW}{\mathcal{W}}
\newcommand{\cX}{\mathcal{X}}  \newcommand{\cY}{\mathcal{Y}}
\newcommand{\cZ}{\mathcal{Z}}

% --- RL-specific notation ---
\newcommand{\states}{\mathcal{S}}
\newcommand{\actions}{\mathcal{A}}
\newcommand{\policy}{\pi}
\newcommand{\Vpi}{V^{\pi}}
\newcommand{\Qpi}{Q^{\pi}}
\newcommand{\Vstar}{V^{*}}
\newcommand{\Qstar}{Q^{*}}
\newcommand{\Bell}{\mathcal{T}}
\newcommand{\Bellpi}{\mathcal{T}^{\pi}}
\newcommand{\disc}{\gamma}

% --- Citation helpers ---
\newcommand\citb[1]{{\texttt{\textcolor{bleucit}{[#1]}}}}
\newcommand\citbs[1]{{\small\texttt{\textcolor{bleucit}{[#1]}}}}
\newcommand\citw[1]{{\texttt{\textcolor{white}{[#1]}}}}

% --- Marks ---
\newcommand{\cmark}{{\color{MyGreen}\ding{51}}}
\newcommand{\xmark}{{\color{MyRed}\ding{55}}}
\newcommand{\mybullet}{{\color{MyGreen}\bullet}}

% --- Figure path ---
\graphicspath{{Fig/}}

% ============================================================
% TITLE — update for each week
% ============================================================
\title{Week XX — Topic Title}
\author[Andreu Cecilia]{\large \textbf{Andreu Cecilia}}
\date{240MAR45 — Reinforcement Learning in Robotics \\ 2026}

% ============================================================
\begin{document}
% ============================================================

\begin{frame}
  \titlepage
  \begin{figure}
    \centering
    \includegraphics[width=0.3\textwidth]{Fig/cropped-logo-ETSEIB-UPC}
  \end{figure}
\end{frame}

% ============================================================
\section{Section Title}
% ============================================================

\begin{frame}{Frame Title}

  \blueblock{Key idea}{
    Main conceptual point goes here.
  }

  \vspace{3mm}

  Supporting content, equation, or figure.

\end{frame}

% ============================================================
% RECURRING ARTIFACTS — copy-paste as needed
% ============================================================

% --- Unifying diagram frame ---
\begin{comment}
\begin{frame}{The Unifying Diagram}
  \begin{center}
  \begin{tabular}{l|cc}
    \toprule
                      & \textbf{Model known} & \textbf{Model unknown} \\
    \midrule
    \textbf{Value-based}  & DP / LQR             & Q-learning, DQN        \\
    \textbf{Policy-based} & iLQR                 & REINFORCE              \\
    \textbf{Hybrid}       & MPC                  & Actor-critic, Dyna     \\
    \bottomrule
  \end{tabular}
  \end{center}
  \blueblock{Today}{Highlight the cell of the method introduced this week.}
\end{frame}
\end{comment}

% --- Failure-modes frame (Week 7+, open every theory session) ---
\begin{comment}
\begin{frame}{Failure Modes of Approximate RL}
  \redblock{Not a complete taxonomy}{
    These are the three failure modes that drive the course narrative.
    Other failure modes exist (non-stationarity, multi-agent, deployment failures).
  }
  \begin{enumerate}
    \item \textbf{Instability} under function approximation \hfill \citbs{Week 7: DQN}
    \item \textbf{High variance} in policy optimization \hfill \citbs{Week 9: baselines, GAE}
    \item \textbf{Distribution shift} in data \hfill \citbs{Week 12: offline RL}
  \end{enumerate}
  \vspace{2mm}
  \textit{Other failure modes (touched, not centered):}
  \begin{enumerate}\setcounter{enumi}{3}
    \item Insufficient exploration \hfill \citbs{Week 4+}
    \item Model bias \hfill \citbs{Week 11}
    \item Reward misspecification \hfill \citbs{Week 13}
  \end{enumerate}
  \blueblock{Today}{\color{MyRed}\textbf{Failure mode X} — [name] — is the focus of today's session.}
\end{frame}
\end{comment}

% --- Actor-critic arc frames (Weeks 7 / 9 / 10) ---
\begin{comment}
% End of Week 7:
\blueblock{Coming up after the exam}{
  Today we covered \textbf{failure mode 1} (instability under FA). \\
  After the exam: \textbf{failure mode 2} (high variance in policy optimization). \\
  Week 10: one construction addresses both. Actor-critic is not a topic --- it is a conclusion.
}

% Start of Week 9:
\blueblock{Recap: where we are}{
  Two weeks ago: value-based methods break under approximation (failure mode 1). \\
  Today: direct policy optimization --- failure mode 2. \\
  Next week: why neither branch alone is sufficient.
}

% End of Week 9:
\blueblock{Hold this thought}{
  Two methods. Two dominant failure modes. Two weaknesses. \\
  Next week: one construction fixes both.
}

% Start of Week 10:
\blueblock{Actor-critic is not a topic --- it is a conclusion}{
  Failure mode 1 makes value-based methods unstable. \\
  Failure mode 2 makes policy-based methods high-variance. \\
  Today: one construction addresses both.
}
\end{comment}

\end{document}
